\section{A square-root threshold for density sublevels}
\label{uni-sec:threshold}

The topology of the positive-density set can be controlled by the volume of small sublevels. A nonnegative $C^{1,1}$ density grows at most quadratically away from a zero. If its zero set separates a small ball, a family of transverse segments therefore contributes at least a constant times $\sqrt t$ to the volume of $\{f\leq t\}$. Excluding that lower bound gives connectedness and hence uniqueness.

Write $\operatorname{Vol}(E)=\int_E\mathrm{d}V$ for Borel sets $E$.

\begin{theorem}\label{uni-thm:threshold}
Let $u,v\in C^2(X)$, $A,B\geq0$, and
$A^n=B^n=f\,\mathrm{d}V$, where $f\geq0$ is $C^{1,1}$.  If
\begin{equation}
  \liminf_{t\downarrow0}
  \frac{\operatorname{Vol}\{f\leq t\}}{\sqrt t}=0,
  \label{uni-eq:sqrtcriterion}
\end{equation}
then $\{f>0\}$ is connected and $u-v$ is constant.  In particular this
holds when $f^{-1/2}\in L^1(\mathrm{d}V)$, with $f^{-1/2}=+\infty$ on $\{f=0\}$.
The exponent $1/2$ is sharp in every dimension $n\geq3$.
\end{theorem}

\begin{proof}
Write $Z=\{f=0\}$.  Condition \eqref{uni-eq:sqrtcriterion} first implies
$\operatorname{Vol}(Z)=0$.  We show that the complement of $Z$ in every sufficiently
small coordinate ball is connected.  Fix such a ball with convex Euclidean
image.  If its complement were disconnected, choose points $a,b$ in
different components and small closed endpoint disks perpendicular to the
segment $[a,b]$, contained in those respective components.  By making
the disks small, all parallel segments between corresponding points stay
in the ball.  Each segment meets $Z$, since otherwise it would connect
the two components.

On a slightly smaller fixed coordinate region the gradient of $f$ is
Lipschitz with some constant $K>0$.  Nonnegativity gives
$\nabla f(z)=0$ at every zero, hence
\begin{equation}\label{uni-eq:quadratic-upper}
 f(x)\leq\frac K2|x-z|^2
\end{equation}
whenever the segment from $z$ to $x$ stays in that region.  The endpoint
disks are compact and disjoint from $Z$, so their nearby segment portions
are uniformly free of zeros.  Thus every parallel segment contains a zero
a fixed positive distance from its endpoints.  For small $t>0$,
\eqref{uni-eq:quadratic-upper} supplies on that segment an interval of length
$2\sqrt{2t/K}$ belonging to $\{f\leq t\}$.  Integration over the
positive-measure transverse disk, and comparability of $\mathrm{d}V$ with
coordinate volume, give
\[
 \operatorname{Vol}\{f\leq t\}\geq c\sqrt t
\]
for all sufficiently small $t$, contradicting
\eqref{uni-eq:sqrtcriterion}.  This proves the local connectedness assertion.

Each such ball has a nonempty complement of $Z$, since $Z$ has zero
volume, and that complement belongs to one component of $U=X\setminus Z$.
Two overlapping balls have the same component label: their open overlap
contains a point outside $Z$.  These locally constant labels on
connected $X$ are therefore all the same.  Hence $U$ is connected,
and Theorem~\ref{uni-thm:anchors} gives uniqueness.

If $f^{-1/2}\in L^1$, then, outside the null zero set,
\[
 \frac{\operatorname{Vol}\{f\leq t\}}{\sqrt t}
 \leq\int_{\{f\leq t\}}f^{-1/2}\,\mathrm{d}V\longrightarrow0
\]
by dominated convergence.  This proves the stated sufficient integrability
condition.  The sharpness assertion follows from the examples in
Section~\ref{uni-sec:examples}, as quantified at the end of that section.
\end{proof}

\begin{remark}
The regularity in this criterion is a condition on the \emph{density}:
the quadratic upper bound at its zeros is the decisive input.
It is different from the $C^{1,1}$ condition on the \emph{potentials}
in the preceding uniqueness theorem.
\end{remark}

